% !TEX root = ../../ctfp-print.tex

\lettrine[lhang=0.17]{自}{由构造是}伴随的一个强有力的应用. \newterm{free functor 自由函子} 被定义
为 \newterm{forgetful functor 遗忘函子} 的左伴随. 遗忘函子通常是个相当简单的函子, 它忘掉某些结构.
举例来说, 许多有趣的范畴是搭建在集合之上的. 但抽象掉这些集合之后的范畴物件没有任何内部结构 ---
它们没有元素. 话虽如此, 这些物件往往仍带着关于集合的记忆: 存在一个从给定范畴 $\cat{C}$ 到
$\Set$ 的映射 --- 一个函子. 与 $\cat{C}$ 中某个物件相对应的那个集合, 叫作它的
\newterm{underlying 底层} 集合.

幺半群就是这类拥有底层集合的物件 --- 元素的集合. 存在一个遗忘函子 $U$, 从幺半群的范畴
$\cat{Mon}$ 到集合的范畴, 它把幺半群映到它们的底层集合. 它还把幺半群的态射 (同态) 映成集合之间
的函数.

我喜欢把 $\cat{Mon}$ 想象成有着双重人格. 一方面, 它是一堆带着乘法和单位元的集合. 另一方面, 它是
一个范畴, 其中的物件毫无特征, 唯一的结构都编码在它们彼此之间的态射里. 每个保持乘法与单位元的
集合函数, 都导出 $\cat{Mon}$ 中的一个态射.\\
\newline
需要记在心里的几件事:

\begin{itemize}
  \tightlist
  \item
        可能有多个幺半群映到同一个集合, 而且
  \item
        幺半群的态射, 比它们底层集合之间的函数更少 (或至多一样多).
\end{itemize}

\noindent
那个作为遗忘函子 $U$ 之左伴随的函子 $F$, 就是从生成集搭建自由幺半群的自由函子. 这个伴随
源于我们先前讨论过的自由幺半群的普遍构造.\footnote{参见第 13 章关于
  \hyperref[free-monoids]{自由幺半群}的内容.}

\begin{figure}[H]
  \centering
  \includegraphics[width=0.6\textwidth]{images/forgetful.jpg}
  \caption{幺半群 $m_1$ 和 $m_2$ 有着相同的底层集合. $m_2$ 与 $m_3$ 的底层集合之间的
    函数, 比它们之间的态射更多.}
\end{figure}

\noindent
用 hom-集来表达, 我们可以把这个伴随写成:
\[\cat{Mon}(F x, m) \cong \Set(x, U m)\]
这个同构 (对 $x$ 和 $m$ 都是自然的) 告诉我们:

\begin{itemize}
  \tightlist
  \item
        对由 $x$ 生成的自由幺半群 $F x$ 与任意幺半群 $m$ 之间的每个幺半群同态, 都存在唯一
        的函数, 把生成元的集合 $x$ 嵌入 $m$ 的底层集合. 它是 $\Set(x, U m)$ 中的一个函数.
  \item
        对把 $x$ 嵌入某个 $m$ 的底层集合的每个函数, 都存在唯一的幺半群态射, 从由 $x$ 生成的
        自由幺半群走到幺半群 $m$. (这个态射就是我们在普遍构造里叫作 $h$ 的那个.)
\end{itemize}

\begin{figure}[H]
  \centering
  \includegraphics[width=0.8\textwidth]{images/freemonadjunction.jpg}
\end{figure}

\noindent
直觉是, $F x$ 是以 $x$ 为基础所能搭建的 ``最大'' 幺半群. 如果我们能看透幺半群的内部, 就会
看到: 属于 $\cat{Mon}(F x, m)$ 的任何态射都把 \emph{这个} 自由幺半群\emph{嵌入}了另一个
幺半群 $m$. 它做到这一点的方式, 是把某些元素互相认同 (也可能一个都不认同). 特别地, 它把
$F x$ 的生成元 (也就是 $x$ 的元素) 嵌入 $m$. 伴随表明: 右边那个由 $\Set(x, U m)$ 中的函数
给出的 $x$ 的嵌入, 唯一地决定了左边幺半群之间的嵌入, 反之亦然.

在 Haskell 里, 列表这个数据结构就是一个自由幺半群 (有几点需要留意: 参见
\urlref{http://comonad.com/reader/2015/free-monoids-in-haskell/}{Dan Doel 的博客文章}).
列表类型 \code{{[}a{]}} 是一个自由幺半群, 其中类型 \code{a} 就代表生成元的集合. 举例来说,
类型 \code{{[}Char{]}} 包含单位元 --- 空列表 \code{{[}{]}} --- 以及像 \code{{[}'a'{]}},
\code{{[}'b'{]}} 这样的单例 --- 它们就是这个自由幺半群的生成元. 其余的东西都由施加
``乘法'' 生成. 在这里, 两个列表的乘法不过是把一个接到另一个后面. 接合满足结合律, 而且有单位元
(也就是说存在一个中立元素 --- 这里是空列表). 由 \code{Char} 生成的自由幺半群, 无非就是所有由
\code{Char} 中的字符组成的字符串的集合. 它在 Haskell 里叫 \code{String}:

\src{snippet01}
(\code{type} 定义的是类型别名 --- 给一个已有的类型另起一个名字.)

另一个有趣的例子, 是只由一个生成元搭建的自由幺半群. 它就是单位值的列表类型
\code{{[}(){]}}. 它的元素是 \code{{[}{]}}, \code{{[}(){]}}, \code{{[}(), (){]}},
等等. 每个这样的列表都可以由一个自然数 --- 它的长度 --- 来描述. 单位值的列表里没有编码更多的
信息. 接合这样两个列表, 产生的新列表的长度是两者长度之和. 很容易看出, 类型 \code{{[}(){]}}
同构于自然数的加法幺半群 (含零). 这里是一对互为逆的函数, 见证了这个同构:

\src{snippet02}
为了简单起见, 我用了类型 \code{Int} 而不是 \code{Natural}, 但想法是一样的. 函数
\code{replicate} 创建一个长度为 \code{n}、预先填满某个给定值的列表 --- 这里填的是单位元.

\section{一些直觉}

接下来是一些手挥 (hand-waving) 式的论证. 这类论证远谈不上严格, 但它们有助于形成直觉.

要对自由/遗忘伴随建立直觉, 记住这一点会有帮助: 函子和函数在本性上都是有损的. 函子可能把多个
物件、多个态射折叠到一处, 函数可能把集合的多个元素归成一堆. 此外, 它们的像也可能只覆盖陪域的
一部分.

$\Set$ 里一个 ``平均'' 水平的 hom-集, 会容纳一整个谱系的函数: 从有损程度最低的 (比如单射,
或者干脆就是同构) 开始, 一路到把整个定义域折叠成单个元素的常数函数 (如果定义域里有元素的
话).

我倾向于把任意范畴里的态射也想成是有损的. 这不过是个心智模型, 但它很有用, 尤其在思考伴随的
时候 --- 特别是其中一个范畴是 $\Set$ 的那些伴随.

严格地说, 我们只能谈可逆的态射 (同构) 与不可逆的态射. 可以被想成有损的是后一类. 还有
单态射与满态射 (mono / epi) 的概念, 它们推广了单射 (不折叠) 与满射 (覆盖整个陪域) 这两个关于
函数的想法; 但完全可能存在一个态射既单又满, 却依然不可逆.

在 自由 ⊣ 遗忘 伴随里, 左边是约束更多的范畴 $\cat{C}$, 右边是约束更少的范畴 $\cat{D}$.
$\cat{C}$ 中的态射 ``更少'', 因为它们必须保持某种额外的结构. 就 $\cat{Mon}$ 而言, 它们必须
保持乘法与单位元. $\cat{D}$ 中的态射不必保持那么多结构, 所以它们 ``更多''.

当我们把遗忘函子 $U$ 作用在 $\cat{C}$ 中的物件 $c$ 上时, 我们把它想成是在揭示 $c$ 的
``内部结构''. 事实上, 若 $\cat{D}$ 就是 $\Set$, 我们会把 $U$ 想成是在\emph{定义} $c$ 的内部
结构 --- 它的底层集合. (在任意范畴里, 除了通过物件彼此之间的连接, 我们无从谈起物件的内部,
但这里我们只是在做手挥论证.)

如果我们用 $U$ 映射两个物件 $c'$ 和 $c$, 那么一般而言, 我们预期 hom-集 $\cat{C}(c', c)$
的像只覆盖 $\cat{D}(U c', U c)$ 的一个子集. 这是因为 $\cat{C}(c', c)$ 中的态射必须保持那份额
外的结构, 而 $\cat{D}(U c', U c)$ 中的不必.

\begin{figure}[H]
  \centering
  \includegraphics[width=0.45\textwidth]{images/forgettingmorphisms.jpg}
\end{figure}

\noindent
但由于伴随是被定义为特定 hom-集之间的 \newterm{isomorphism 同构}, 我们就得极挑剔地挑选
$c'$. 在伴随里, $c'$ 不是从 $\cat{C}$ 里随便哪儿挑出来的, 而是从自由函子 $F$ 的那个 (大概
更小的) 像里挑出来的:
\[\cat{C}(F d, c) \cong \cat{D}(d, U c)\]
因此 $F$ 的像必须由这样一些物件组成: 有大量态射从它们出发, 走向任意的 $c$. 事实上, 从
$F d$ 到 $c$ 的保持结构的态射, 必须和从 $d$ 到 $U c$ 的不保持结构的态射一样多. 这意味着 $F$
的像基本上由无结构的物件组成 (于是也没什么结构要求态射去保持). 这样的 ``无结构'' 物件就叫作
自由对象 (free objects).

\begin{figure}[H]
  \centering
  \includegraphics[width=0.45\textwidth]{images/freeimage.jpg}
\end{figure}

\noindent
在幺半群这个例子里, 自由幺半群除了由单位元律和结合律生成出来的东西之外, 没有别的结构. 除此
之外, 所有乘法都产出全新的元素.

在一个自由幺半群里, $2 * 3$ 并不是 $6$ --- 它是一个新元素 ${[}2, 3{]}$. 既然
${[}2, 3{]}$ 与 $6$ 之间没有互相认同, 那么这个自由幺半群到任意别的幺半群 $m$ 的态射,
就允许把它们分别映射. 但把 ${[}2, 3{]}$ 和 $6$ (它们的乘积) 映到 $m$ 中同一个元素也没问题.
或者在一个加法幺半群里把 ${[}2, 3{]}$ 与 $5$ (它们的和) 认同起来, 诸如此类. 不同的认同给出
不同的幺半群.

这又引出另一个有趣的直觉: 自由幺半群不去执行幺半群运算, 而是把传进来的参数累积起来. 它们不把
$2$ 和 $3$ 乘起来, 而是把 $2$ 和 $3$ 记在一个列表里. 这套做法的好处是, 我们不必事先指定将要
使用哪种幺半群运算. 我们可以一直累积参数, 直到最后才对结果施加一个运算符. 而也正是在那一刻,
我们才选择要施加哪个运算符. 我们可以把这些数加起来, 或者乘起来, 或者模 2 相加, 如此等等. 一个
自由幺半群把表达式的创建与表达式的求值分离开来. 等我们讲代数的时候还会再见到这个想法.

这个直觉可以推广到其他更精细的自由构造. 比如说, 我们可以先累积起一棵又一棵的表达式树, 再去
求值. 这么做的好处是, 我们可以变换这些树, 让求值更快或者更省内存. 实现矩阵演算时就是这样做的
--- 若采用急切求值 (eager evaluation), 就会为了存放中间结果而大量分配临时数组.

\section{挑战}

\begin{enumerate}
  \tightlist
  \item
        考虑一个以单例集合为生成元搭建起来的自由幺半群. 证明: 从这个自由幺半群到任意幺半群
        $m$ 的态射, 与从那个单例集合到 $m$ 的底层集合的函数之间, 存在一一对应.
\end{enumerate}
